%#!PATH=/opt/texlive/2025/bin/universal-darwin:$PATH; ptex2pdf -l -u -ot '--synctex=1 --shell-escape' -od '-p a4' template-e; skim_reload_pdf template-e
% v1.4 [2026/03/02]
\documentclass[english]{jjsce}
%\documentclass[english,jscefinal]{jjsce}
\usepackage{amsmath}
%\usepackage{amsthm}
\usepackage[defaultsups]{newtxtext}
\usepackage[varg]{newtxmath}
%\usepackage{bm}
%%
\usepackage[dvipdfmx]{graphicx}% (u)pLaTeX + dvipdfmx
%% When using LuaLaTeX, select the following:
%\usepackage{graphicx}
%\usepackage{stfloats}\fnbelowfloat
%%
%\usepackage{xcolor}
\usepackage[superscript]{cite}
\usepackage{url}
\usepackage{endnotes}
\usepackage{enumitem}
%% If you use the hyperref package, select either of the following
%\usepackage[dvipdfmx]{hyperref}% (u)pLaTeX + dvipdfmx
%\usepackage{pxjahyper}% (u)pLaTeX + dvipdfmx
%\usepackage[luatex,pdfencoding=auto]{hyperref}% LuaLaTeX

\begin{document}
\title{INSTRUCTIONS FOR MAUSCRIPT FOR\\
JOURNAL OF JSCE/ AUTHOR GUIDELINES FOR\\
MANUSCRIPT PREPARATION FOR JSCE}
%\subtitle{}
\authorlist{%
 \authorentry{Taro DOBOKU}{doboku-u}
 \authorentry{Hanako YOTSUYA}{doboku-co}
 \authorentry{John SMITH}{u-tokyo}
}
\affiliate[doboku-u]{Member of JSCE, Professor, Dept. of Civil Eng., University of Doboku \\
(Yotsuya 1, Shinjuku-ku, Tokyo 160-0004, Japan)}{doboku@jsce.ac.jp}
\Caffiliate[doboku-co]{Member of JSCE, Dept. of Civil Eng., Doboku Corporation\\
(13-5, Mitsuya 6, Shinjuku-ku, Tokyo 160-0004, Japan) }
{hanako@jsce.co.jp}
\affiliate[u-tokyo]{Member of JSCE, Professor, Inst., Industrial Science, University of Tokyo\\
 (7-22-1 Roppongi, Minato-ku, Tokyo 106-8558. Japan)}{}%{smith@jsce.or.jp}

%\breakauthorline{4}
%\received{2026}{2}{16}% Received Date
%\accepted{2026}{3}{2}% Accepted Date

\begin{abstract}
This template is prepared for your preparion of manscript for JSCE journals. It provides instructions:
page layout, font style, size and others. You may use it to create your own manuscript by replacing the
relevant text with your own, using ``cut \& paste.'' The Abstract should be justified, leaving a 30 mm margine
on the left and right sides. Font should be a 10-point Times-New-Roman. The length should be 300 words
or less.
\end{abstract}
\begin{keyword}
times, italic, 10pt, one blank line below abstract, indent if key words exceed one line 
\end{keyword}
\maketitle

\section{Title PAGE}

The first page consists of two parts:
%\begin{description}
% \item[(a)] Front matters : single column (title, author(s),
%affiliation(s), contact address(es), E-mail address(es), abstract, key words)．The e-mail address of the corresponding author is mandatory.
%E-mail address should be indicated in a separate,
%independent line.
% \item[(b)] Main text in double columns.
%\end{description}
\begin{enumerate}[label=(\alph*)]
 \item Front matters : single column (title, author(s),
affiliation(s), contact address(es), E-mail address(es), abstract, key words)．The e-mail address of the corresponding author is mandatory.
E-mail address should be indicated in a separate,
independent line.\\[-20pt]
 \item Main text in double columns.\\[-20pt]
\end{enumerate}

The journal name, volume and issue numbers and
the date of issue should be aligned right in the top
margin. Page numbers are to be put in the bottom
margins of the manuscript. Some word processing
softwares do not allow texts in both single and double
columns to be put together in one file, and thus, create
two separate files for the title page.

\subsection{Layout and fonts for the front matters}

Left and right margins for the front matters are
equally set at 30 mm, The front matters are, thus, to
be laid-out within the borders narrower than those for
the main text.

The front matters include the followings:\\
\hspace*{40pt}(About 10 mm blank space)

\textbf{Title} in Times-New-Roman, 18pt, bold

\hspace*{40pt}(About 15mm blank space)

\textbf{Author(s)} in Times-New-Roman, 12pt.

\hspace*{40pt} (About 5 mm blank space)

\textbf{Affiliation(s)} in Times-New-Roman, 9pt.

\textbf{E-mail address(es)} in Times-New-Roman, 9pt

\hspace*{40pt} (About 10 mm blank space)

\textbf{Abstract} in Times-New-Roman, 10pt, max. 300 words,

\hspace*{40pt} (1 line spacing) and

\textbf{About 5 Key Words} in Times-Italic, 10pt,

\hspace*{40pt} (max. 2 blank lines).

The name(s) and affiliation(s) of the author(s)
should be numbered in order of appearance as shown
above. The title \textbf{'\textit{Key Words}'} is bold and italic. 

\subsection{Layout and fonts of the main text}

The text should be placed about 1cm below the key words.
Left and right margins for the text are equally set at 20 mm.
The text, in double columns put side by
side with 6 mm gap in between, must be singlespaced with
 double spacing between chapters. 
Use 11pt Times-New-Roman font for the text. 

\subsection{page number}

Page numbers should be center-aligned and should
appear at the bottom of each page. 
Since these pieces of information will be notified by the secretariat of
JSCE before completing the final manuscripts, leave
the blanks as they are, and number the pages tentatively from 1. 

%\Acknowledgment

%\appendix
%\section{付録の見出し}

\begin{thebibliography}{9}
\bibitem{1}
...
\end{thebibliography}

\end{document}

\begin{table}[htb]
\caption{}
%\label{}
\centering
\begin{tabular}{ll}
\hline
%
\hline
\end{tabular}
\end{table}

\begin{figure}[htb]
\centering
% graphic
\caption{}
%\label{}
\end{figure}
